\section{Materials and Methods}

\subsection{EEG Hardware: Bitbrain Air 8-Channel System}

The choice of EEG hardware in clinical populations must balance signal quality, patient comfort, and practical usability. For this proof-of-concept study, we selected the Bitbrain Air wireless EEG system with 8 dry electrodes, specifically designed for ease of use in real-world clinical environments.

\begin{table}[H]
\centering
\caption{Technical specifications of the Bitbrain Air EEG system}
\label{tab:hardware}
\begin{tabular}{@{}ll@{}}
\toprule
\textbf{Parameter} & \textbf{Specification} \\
\midrule
Number of channels & 8 \\
Electrode type & Dry Ag/AgCl active electrodes \\
Montage & Reduced 10-20 system (F3, F4, C3, C4, P3, P4, O1, O2) \\
Sampling rate & 256 Hz \\
ADC resolution & 24-bit \\
Dynamic range & $\pm$187.5 mV \\
Wireless transmission & Bluetooth 4.0 (range up to 10 m) \\
Battery life & 8 hours continuous recording \\
Weight & 90 grams \\
Impedance monitoring & Automatic, < 50 k$\Omega$ target \\
\bottomrule
\end{tabular}
\end{table}

The system offers several key advantages for clinical application:

\begin{itemize}
    \item \textbf{No gel}: Eliminates lengthy preparation and cleanup (2-3 minutes setup)
    \item \textbf{Comfort}: Lightweight design reduces skin irritation and discomfort
    \item \textbf{Signal quality}: Active electrodes provide SNR comparable to wet systems up to 40 Hz
    \item \textbf{Mobility}: Wireless design allows recording during natural care activities
    \item \textbf{Real-time access}: Python SDK provides low-latency data streaming (< 50 ms)
\end{itemize}

While 8 channels provide less spatial information than high-density systems, they represent an optimal trade-off for this clinical context, providing sufficient coverage for frontal asymmetry, central motor/arousal regions, temporal social-emotional areas, and posterior attentional regions.

\subsection{Participant and Recording Protocol}

For technical validation, we recruited a healthy female volunteer aged 42 years from the laboratory personnel. The participant had no history of neurological or psychiatric disorders and provided informed consent. The session lasted approximately 5 minutes and was conducted in a quiet room with the participant seated comfortably.

During the session, the participant was instructed to simulate different emotional states based on standardized verbal cues provided by the experimenter every 30-60 seconds. The simulated states included: fatigue, thirst (soif), serenity (sérénité), desire (envie), joy (joie), and neutral. A researcher provided real-time annotations through the tagging interface, resulting in 9 annotated events.

\begin{table}[H]
\centering
\caption{Participant characteristics}
\label{tab:participant}
\begin{tabular}{@{}lc@{}}
\toprule
\textbf{Characteristic} & \textbf{Value} \\
\midrule
Subject ID & F42 \\
Age (years) & 42 \\
Sex & Female \\
Session duration (seconds) & 283.9 \\
Number of annotations & 9 \\
Unique emotional states & 6 \\
\bottomrule
\end{tabular}
\end{table}

\subsection{Real-Time Acquisition Pipeline}

The Bitbrain Python SDK was used to stream data in real time. A custom acquisition script, \texttt{01\_acquisition.py}, was developed with the following features:

\begin{itemize}
    \item Low-latency streaming of EEG blocks with timestamp synchronization
    \item Circular buffering to prevent packet loss during high CPU load
    \item Real-time impedance monitoring and display
    \item Optional computation of band powers and frontal alpha asymmetry for visualization
    \item Continuous logging of raw EEG into timestamped CSV files
\end{itemize}

The raw EEG log format is:
\begin{verbatim}
time_s,ch1,ch2,ch3,ch4,ch5,ch6,ch7,ch8
0.000,12.34,11.56,13.21,...
\end{verbatim}

\subsection{Real-Time Emotional Annotation Interface}

Because patients cannot self-report, emotional labeling must rely on real-time observation by caregivers. We developed a simple graphical interface using Streamlit (\texttt{02\_tagging\_app.py}) with the following features:

\begin{itemize}
    \item "Start Session" button defining time zero and synchronizing with EEG
    \item Large, color-coded buttons for each observable state
    \item Text input for contextual notes
    \item Real-time display of recent annotations
    \item Session management with unique session IDs
\end{itemize}

Each button click appends a line to \texttt{annotations\_patient.csv}:
\begin{verbatim}
time_s,label,session_id,notes
6.962,fatigue,20260304_020920,
10.158,soif,20260304_020920,
14.522,serenite,20260304_020920,
\end{verbatim}

\subsection{EEG Preprocessing}

Raw EEG data were preprocessed using MNE-Python version 1.6.0 with the following pipeline:

\begin{enumerate}
    \item \textbf{Band-pass filtering}: 1 to 40 Hz using a zero-phase FIR filter to remove slow drifts and high-frequency noise
    \item \textbf{Notch filtering}: 50 Hz filter to remove power line interference
    \item \textbf{Artifact detection}: Channel-wise amplitude thresholds of $\pm$150 $\mu$V and gradient thresholds of 50 $\mu$V/ms
    \item \textbf{Independent Component Analysis (ICA)}: Infomax ICA to identify and remove components corresponding to eye movements, blinks, and muscle artifacts
    \item \textbf{Channel interpolation}: Bad channels, if any, were interpolated using spherical spline interpolation
    \item \textbf{Re-referencing}: Data were re-referenced to the common average
\end{enumerate}

After preprocessing, data quality metrics were computed including percentage of retained data, signal-to-noise ratio (SNR), and number of rejected ICA components.

\subsection{EEG Windowing and Feature Extraction}

Continuous EEG was segmented into overlapping windows of duration 4 seconds with a stride of 1 second (75\% overlap). This window size was chosen based on optimization studies showing optimal trade-off between temporal resolution and spectral estimation stability \cite{dhuli2022tutorial}. For each window, we extracted a comprehensive feature vector.

\subsubsection{Spectral Features}

Power spectral density (PSD) was computed using Welch's method with 1-second Hann windows and 50\% overlap, resulting in a frequency resolution of 1 Hz. Band powers were extracted for:
\begin{itemize}
    \item Delta (1-4 Hz)
    \item Theta (4-8 Hz)
    \item Alpha (8-13 Hz)
    \item Beta (13-30 Hz)
\end{itemize}

Band powers were normalized by total power in the 1-40 Hz range to reduce inter-session variability:
\begin{equation}
P_{\text{norm}}(b) = \frac{P(b)}{\sum_{b'} P(b')}
\end{equation}

\subsubsection{Asymmetry Features}

Frontal alpha asymmetry (FAA) was computed as:
\begin{equation}
\text{FAA} = \log(\alpha_{\text{right}}) - \log(\alpha_{\text{left}})
\end{equation}
where $\alpha_{\text{right}}$ and $\alpha_{\text{left}}$ represent average alpha power over right frontal (F4) and left frontal (F3) channels respectively.

Regional alpha powers were also extracted for left and right frontal areas individually.

\subsubsection{Connectivity Features}

Functional connectivity was assessed using Pearson correlation coefficients between all channel pairs:
\begin{equation}
r_{ij} = \frac{\text{cov}(x_i, x_j)}{\sigma_i \sigma_j}
\end{equation}

From the correlation matrix, we extracted two summary statistics:
\begin{itemize}
    \item \textbf{Mean absolute correlation}: $\frac{1}{N} \sum_{i<j} |r_{ij}|$
    \item \textbf{Standard deviation of absolute correlation}: $\sqrt{\frac{1}{N-1} \sum_{i<j} (|r_{ij}| - \bar{r})^2}$
\end{itemize}

These metrics capture overall connectivity strength and its variability across channel pairs.

\subsubsection{Complete Feature Set}

The complete feature vector for each window thus consisted of 9 dimensions:
\begin{enumerate}
    \item $\text{BP}_{\delta}$: Delta band power (normalized)
    \item $\text{BP}_{\theta}$: Theta band power (normalized)
    \item $\text{BP}_{\alpha}$: Alpha band power (normalized)
    \item $\text{BP}_{\beta}$: Beta band power (normalized)
    \item $\alpha_{\text{left}}$: Left frontal alpha power
    \item $\alpha_{\text{right}}$: Right frontal alpha power
    \item $\text{FAA}$: Frontal alpha asymmetry
    \item $\text{mean\_abs\_corr}$: Mean absolute correlation
    \item $\text{std\_abs\_corr}$: Standard deviation of absolute correlation
\end{enumerate}

All features were normalized per subject using z-scoring:
\begin{equation}
z = \frac{x - \mu}{\sigma}
\end{equation}
where $\mu$ and $\sigma$ are the mean and standard deviation across all windows for that subject.

\subsection{Self-Organizing Map (SOM) Modeling}

For each subject, we trained a two-dimensional Self-Organizing Map to obtain a low-dimensional, topology-preserving representation of EEG feature space.

\subsubsection{SOM Architecture}

\begin{itemize}
    \item \textbf{Grid size}: 10 $\times$ 10 hexagonal grid (100 neurons)
    \item \textbf{Input dimension}: 9 (feature vector size)
    \item \textbf{Distance metric}: Euclidean distance in feature space
    \item \textbf{Neighborhood function}: Gaussian with initial radius 3, decreasing exponentially
    \item \textbf{Learning rate}: Initial 0.5, decreasing linearly to 0.01
    \item \textbf{Training epochs}: 500 with batch training algorithm
    \item \textbf{Initialization}: Linear initialization along principal eigenvectors
\end{itemize}

\subsubsection{SOM Training Algorithm}

For each input vector $\mathbf{x}$, the best matching unit (BMU) $\mathbf{m}_c$ is found by:
\begin{equation}
\|\mathbf{x} - \mathbf{m}_c\| = \min_i \|\mathbf{x} - \mathbf{m}_i\|
\end{equation}

The weight vectors are updated according to:
\begin{equation}
\mathbf{m}_i(t+1) = \mathbf{m}_i(t) + \alpha(t) h_{ci}(t) [\mathbf{x}(t) - \mathbf{m}_i(t)]
\end{equation}
where $\alpha(t)$ is the learning rate and $h_{ci}(t)$ is the neighborhood function:
\begin{equation}
h_{ci}(t) = \exp\left(-\frac{\|\mathbf{r}_c - \mathbf{r}_i\|^2}{2\sigma^2(t)}\right)
\end{equation}

\subsubsection{SOM Quality Assessment}

SOM quality was assessed using three metrics:

\begin{enumerate}
    \item \textbf{Quantization Error (QE)}: Average distance between each feature vector and its BMU:
    \begin{equation}
    \text{QE} = \frac{1}{N} \sum_{i=1}^N \|\mathbf{x}_i - \mathbf{m}_{c_i}\|
    \end{equation}
    Lower QE indicates better map resolution.
    
    \item \textbf{Topographic Error (TE)}: Proportion of vectors for which the first and second BMUs are not adjacent:
    \begin{equation}
    \text{TE} = \frac{1}{N} \sum_{i=1}^N u(\mathbf{x}_i)
    \end{equation}
    where $u(\mathbf{x}_i) = 1$ if first and second BMUs are non-adjacent, 0 otherwise. Lower TE indicates better topology preservation.
    
    \item \textbf{Silhouette Coefficient}: For labeled windows, measures cluster separation:
    \begin{equation}
    s(i) = \frac{b(i) - a(i)}{\max\{a(i), b(i)\}}
    \end{equation}
    where $a(i)$ is mean distance to other windows in same cluster and $b(i)$ is mean distance to nearest other cluster.
\end{enumerate}

\subsection{Annotation-EEG Fusion}

To merge SOM trajectories with annotations, we implemented a script that translates annotation events into intervals. Each annotation event defines a candidate emotional state from its time until the next annotation event. Each EEG window with center time inherits the label of the interval containing that time. Windows within 2 seconds of an annotation boundary were excluded to avoid transition effects.

\subsection{Supervised Classification}

We trained a Random Forest classifier with the following parameters:
\begin{itemize}
    \item Number of trees: 200
    \item Maximum depth: None (fully grown)
    \item Minimum samples split: 2
    \item Minimum samples leaf: 1
    \item Class weight: Balanced
    \item Random seed: 42
\end{itemize}

Due to the limited number of annotated windows (n=9), we employed leave-one-out cross-validation to maximize data usage while obtaining unbiased performance estimates. For each fold, 8 windows were used for training and 1 for testing.

Performance metrics included:
\begin{itemize}
    \item Accuracy: $\frac{TP + TN}{TP + TN + FP + FN}$
    \item Precision: $\frac{TP}{TP + FP}$
    \item Recall: $\frac{TP}{TP + FN}$
    \item F1-score: $2 \times \frac{\text{precision} \times \text{recall}}{\text{precision} + \text{recall}}$
\end{itemize}

\subsection{Feature Importance Analysis}

Random Forest feature importance was computed using the Gini impurity decrease:
\begin{equation}
\text{Importance}(f) = \frac{1}{T} \sum_{t=1}^T \sum_{n \in t: \text{split on } f} \Delta G(n)
\end{equation}
where $\Delta G(n)$ is the decrease in Gini impurity at node $n$ when splitting on feature $f$, and $T$ is the number of trees.

To assess stability, we performed bootstrap resampling (1000 iterations) to compute 95\% confidence intervals for each feature's importance score.

\subsection{Connectivity and Spectral Analyses}

For each emotional state, we extracted all 4-second windows centered on annotations and computed:
\begin{itemize}
    \item \textbf{Average connectivity matrix}: Mean Pearson correlation across windows for that state
    \item \textbf{Average power spectrum}: Mean PSD across channels and windows using Welch's method
\end{itemize}

Statistical differences between states were assessed using permutation tests (10,000 permutations) with false discovery rate (FDR) correction for multiple comparisons.

\subsection{Statistical Analysis of Current Results}

For the validation dataset (n=279 windows, 9 annotations), we assessed:

\begin{itemize}
    \item \textbf{Classification significance}: Permutation test with 1000 random label shuffles to establish chance level (16.7\%) and compute p-value
    \item \textbf{Feature importance stability}: Bootstrap resampling (1000 iterations) to compute 95\% confidence intervals for importance scores
    \item \textbf{Cluster separation}: Silhouette coefficient with confidence intervals
    \item \textbf{Connectivity differences}: Paired permutation tests between emotional states with FDR correction
\end{itemize}

All statistical tests were two-tailed with $\alpha = 0.05$.

\subsection{Software and Reproducibility}

To ensure reproducibility:

\begin{itemize}
    \item All analysis code is available at: \url{https://github.com/...}
    \item Random seeds were fixed (seed=42) for all stochastic algorithms
    \item Data processing pipeline is documented in MNE-Python scripts
    \item Parameter choices are justified with citations to literature
\end{itemize}

The following open-source libraries were used:
\begin{itemize}
    \item \textbf{Data processing}: MNE-Python v1.6.0, NumPy v1.21, SciPy v1.7
    \item \textbf{Machine learning}: scikit-learn v1.0, MiniSom v2.3
    \item \textbf{Visualization}: Matplotlib v3.4, Seaborn v0.11
    \item \textbf{Interface}: Streamlit v1.10
\end{itemize}

\subsection{Proposed Clinical Study Design}

Based on the successful validation of the technical pipeline, we have designed a comprehensive clinical study to be conducted pending funding.

\subsubsection{Patient Population}

Eight patients with polyhandicap and expressive aphasia will be recruited from rehabilitation units and long-term care facilities in the Marseille region.

\textbf{Inclusion criteria:}
\begin{itemize}
    \item Diagnosis of polyhandicap with severe expressive aphasia
    \item Preserved vigilance sufficient for EEG recording (>5 minutes eyes open)
    \item Availability of primary caregiver during sessions
    \item Stable medical condition with no acute illness
    \item Age > 18 years
\end{itemize}

\textbf{Exclusion criteria:}
\begin{itemize}
    \item Recent changes in anti-epileptic medication (<3 months)
    \item Open scalp wounds or skin conditions
    \item Active status epilepticus
    \item Refusal by legal guardian
\end{itemize}

\subsubsection{Data Acquisition Protocol}

Each patient will complete up to six recording sessions over a three-month period:
\begin{itemize}
    \item Session duration: 60-90 minutes
    \item Usable recording time: approximately 45 minutes per session
    \item Total windows per patient: approximately 2,700 (45 min $\times$ 60 s/min $\times$ 1 window/s)
    \item Expected annotations: 300-500 per patient
\end{itemize}

Sessions will be scheduled during routine care activities:
\begin{itemize}
    \item Morning care (bathing, dressing)
    \item Feeding times
    \item Physical therapy sessions
    \item Family visits
    \item Quiet rest periods
\end{itemize}

\subsubsection{Caregiver Training}

Primary caregivers will receive standardized training including:
\begin{itemize}
    \item Recognition of behavioral indicators for each emotional category (2 hours)
    \item Use of the annotation interface (30 minutes)
    \item Importance of timely annotation (<3 seconds)
    \item When to use "uncertain" rather than forcing a label
    \item Reference card with behavioral descriptors provided at each session
\end{itemize}

\subsubsection{Outcome Measures}

\textbf{Primary outcomes:}
\begin{itemize}
    \item Data quality and retention rates in patient population
    \item SOM convergence and stability metrics (QE, TE)
    \item Classification accuracy for broad emotional categories
    \item Identification of most discriminative EEG features per patient
\end{itemize}

\textbf{Secondary outcomes:}
\begin{itemize}
    \item Caregiver acceptability (structured interviews)
    \item Correlation with clinical pain and discomfort scales
    \item Preliminary evidence of clinical utility
    \item Longitudinal model stability across sessions
\end{itemize}

\subsubsection{Statistical Analysis Plan}

For each patient:
\begin{itemize}
    \item Cross-validated classification metrics (accuracy, precision, recall, F1)
    \item Chance level established through permutation testing (1,000 random shuffles)
    \item Feature importance correlations with clinical variables
    \item Longitudinal stability assessed by training on early sessions, testing on later sessions
\end{itemize}

For group-level analyses:
\begin{itemize}
    \item Mixed-effects models accounting for repeated measures
    \item Correlations between clinical characteristics and classification performance
    \item Effect size estimation for future larger studies
\end{itemize}